% ============================================================================
% AERIS Paper - Section 5: Experimental Setup
% Target Journal: Applied Intelligence (APIN) - Springer
% ============================================================================

\section{Experimental Setup}
\label{sec:experiments}

This section describes the experimental methodology, including simulation
environment, baseline protocols, evaluation metrics, and statistical
validation procedures.

\subsection{Simulation Environment}

\subsubsection{Platform and Implementation}

Experiments were conducted using a custom Python-based simulator with the
following characteristics:

\begin{itemize}
    \item \textbf{Language}: Python 3.10 with NumPy/SciPy for numerical
    computation
    \item \textbf{Random number generation}: Reproducible via fixed seeds
    \item \textbf{Energy model}: CC2420 TelosB-calibrated (see Section 3.2)
    \item \textbf{Channel model}: Log-normal shadowing with configurable
    parameters
    \item \textbf{Source code}: Available at [repository URL] for
    reproducibility
\end{itemize}

\subsubsection{Network Configuration}

Table~\ref{tab:network_params} summarizes the default network parameters
used across all experiments unless otherwise specified.

\begin{table}[htbp]
\centering
\caption{Default Network Parameters}
\label{tab:network_params}
\begin{tabular}{@{}llc@{}}
\toprule
Parameter & Description & Value \\
\midrule
\multicolumn{3}{l}{\textit{Network Topology}} \\
$N$ & Number of nodes & 50--500 \\
$W \times H$ & Network area & 200m $\times$ 200m \\
Deployment & Node distribution & Uniform random \\
BS position & Base station location & (100, 250) \\
\midrule
\multicolumn{3}{l}{\textit{Energy Parameters}} \\
$E_0$ & Initial energy per node & 0.5 J \\
$E_{elec}$ & Electronics energy & 50 nJ/bit \\
$E_{fs}$ & Free-space amplifier & 10 pJ/bit/m$^2$ \\
$E_{mp}$ & Multi-path amplifier & 0.0013 pJ/bit/m$^4$ \\
$E_{DA}$ & Aggregation energy & 5 nJ/bit \\
$L$ & Packet size & 4000 bits \\
\midrule
\multicolumn{3}{l}{\textit{Protocol Parameters}} \\
$p_{CH}$ & Cluster head probability & 0.1 \\
$k_{skel}$ & Skeleton node count & $\sqrt{k}$ \\
$k_{gw}$ & Gateway count & 2 \\
$L_{max}$ & Gateway load limit & 5 \\
\midrule
\multicolumn{3}{l}{\textit{Channel Parameters}} \\
$\eta$ & Path loss exponent & 2.5 \\
$\sigma$ & Shadowing std. dev. & 4 dB \\
$P_{tx}$ & Transmit power & 0 dBm \\
\bottomrule
\end{tabular}
\end{table}

\subsection{Baseline Protocols}

We compare AERIS against four representative baseline protocols spanning
different design paradigms:

\subsubsection{LEACH (Low-Energy Adaptive Clustering Hierarchy)}

The foundational clustering protocol \cite{heinzelman2000leach} with
probabilistic cluster head selection. Represents $O(1)$ latency approaches
with direct CH-to-BS transmission.

\subsubsection{PEGASIS (Power-Efficient GAthering in Sensor Information Systems)}

Chain-based protocol \cite{lindsey2002pegasis} achieving optimal energy
efficiency through sequential data aggregation. Represents $O(n)$ latency
approaches.

\subsubsection{HEED (Hybrid Energy-Efficient Distributed clustering)}

Multi-hop clustering protocol \cite{younis2004heed} with residual-energy-based
cluster head selection. Represents hybrid approaches balancing energy and
latency.

\subsubsection{TEEN (Threshold-sensitive Energy Efficient sensor Network)}

Event-driven protocol \cite{manjeshwar2001teen} with hard/soft thresholds
for data transmission. Represents reactive protocols for time-critical
applications.

\begin{table}[htbp]
\centering
\caption{Baseline Protocol Characteristics}
\label{tab:baseline_chars}
\begin{tabular}{@{}lcccc@{}}
\toprule
Protocol & Latency & Structure & CH Selection & Year \\
\midrule
LEACH & $O(1)$ & Cluster & Probabilistic & 2000 \\
PEGASIS & $O(n)$ & Chain & Sequential & 2002 \\
HEED & $O(1)$ & Cluster & Energy-based & 2004 \\
TEEN & $O(1)$ & Cluster & Threshold & 2001 \\
\textbf{AERIS} & $O(\log n)$ & Hierarchical & Composite & 2026 \\
\bottomrule
\end{tabular}
\end{table}

\subsection{Evaluation Metrics}

\subsubsection{Primary Metrics}

\begin{enumerate}
    \item \textbf{Packet Delivery Ratio (PDR)}: Fraction of generated packets
    successfully delivered to the base station.
    \begin{equation}
    PDR = \frac{\text{Packets received at BS}}{\text{Packets generated by nodes}}
    \label{eq:pdr}
    \end{equation}

    \item \textbf{Energy Consumption}: Total network energy consumed per round.
    \begin{equation}
    E_{total} = \sum_{i=1}^{N} (E_{tx,i} + E_{rx,i} + E_{agg,i})
    \label{eq:energy}
    \end{equation}

    \item \textbf{End-to-End Latency}: Time from data generation to BS delivery,
    measured in transmission hops.
    \begin{equation}
    Latency = H_{avg} \times T_{hop}
    \label{eq:latency}
    \end{equation}
    where $H_{avg}$ is the average hop count and $T_{hop} = 10$ms per hop
    (IEEE 802.15.4 timing).
\end{enumerate}

\subsubsection{Secondary Metrics}

\begin{enumerate}
    \item \textbf{Network Lifetime}: Number of rounds until first node death
    (FND) or 50\% node death.

    \item \textbf{Energy Balance (Jain's Fairness Index)}:
    \begin{equation}
    J = \frac{(\sum_{i=1}^{N} E_i)^2}{N \cdot \sum_{i=1}^{N} E_i^2}
    \label{eq:jain}
    \end{equation}

    \item \textbf{Topology Dynamic Adaptability (TDA)}: New metric introduced
    in this work, measuring protocol robustness under node churn.
    \begin{equation}
    TDA = \frac{PDR_{churn}}{PDR_{baseline}} \times 100\%
    \label{eq:tda}
    \end{equation}
\end{enumerate}

\subsection{Experimental Scenarios}

\subsubsection{Scenario S1: Baseline Comparison}

Standard comparison under uniform deployment with 50--500 nodes. Evaluates
core protocol performance without environmental perturbations.

\subsubsection{Scenario S2: Scalability Analysis}

Network size varied from 50 to 500 nodes in increments of 50. Evaluates
how protocol performance scales with network density.

\subsubsection{Scenario S3: Dynamic Topology (Node Churn)}

Random node failure rates of 10\%, 20\%, and 30\% applied after network
stabilization. Tests protocol robustness using the TDA metric.

\subsubsection{Scenario S4: Regional Failure}

Concentrated node failures within a 50m radius region, simulating localized
disasters. Tests spatial robustness.

\subsubsection{Scenario S5: Environment Variation}

Channel parameters varied to simulate different environments:
\begin{itemize}
    \item Indoor office: $\eta = 2.0$, $\sigma = 3$dB
    \item Industrial: $\eta = 3.0$, $\sigma = 6$dB
    \item Urban dense: $\eta = 3.5$, $\sigma = 8$dB
\end{itemize}

\subsection{Statistical Methodology}

To ensure rigorous and reproducible results, we employ the following
statistical procedures:

\subsubsection{Sample Size}

Each configuration is executed for \textbf{200 independent runs} with
different random seeds. This sample size provides 95\% confidence intervals
within $\pm$2\% of the mean for typical metric variances.

\subsubsection{Hypothesis Testing}

For pairwise protocol comparisons, we use:
\begin{itemize}
    \item \textbf{Normality test}: Shapiro-Wilk test ($\alpha = 0.05$)
    \item \textbf{Variance homogeneity}: Levene's test
    \item \textbf{Parametric comparison}: Welch's t-test (unequal variances)
    \item \textbf{Non-parametric alternative}: Mann-Whitney U test when
    normality is violated
\end{itemize}

\subsubsection{Multiple Comparison Correction}

When comparing AERIS against multiple baselines, we apply
\textbf{Holm-Bonferroni correction} to control family-wise error rate
at $\alpha = 0.05$.

\subsubsection{Effect Size}

We report \textbf{Cohen's $d$} for practical significance:
\begin{equation}
d = \frac{\bar{x}_1 - \bar{x}_2}{s_{pooled}}
\label{eq:cohens_d}
\end{equation}

Interpretation: $|d| < 0.2$ (negligible), $0.2 \leq |d| < 0.5$ (small),
$0.5 \leq |d| < 0.8$ (medium), $|d| \geq 0.8$ (large).

\subsection{Fair Comparison Guarantees}

To ensure fair comparison across all protocols:

\begin{enumerate}
    \item \textbf{Identical seeds}: All protocols use the same random seed
    for each run, ensuring identical node positions and channel realizations.

    \item \textbf{Unified energy model}: All protocols use the same
    \texttt{ImprovedEnergyModel} class with CC2420 parameters.

    \item \textbf{Consistent channel model}: Same log-normal shadowing
    parameters applied to all protocols.

    \item \textbf{Equal simulation duration}: All experiments run for 200
    rounds unless terminated by network death.

    \item \textbf{No protocol-specific tuning}: Default parameters from
    original publications used for baseline protocols.
\end{enumerate}

\subsection{Reproducibility}

All experimental code, configuration files, and raw result data are
available in the supplementary materials. The repository includes:

\begin{itemize}
    \item Complete Python source code for all protocols
    \item Configuration files for all experimental scenarios
    \item Scripts to regenerate all figures and tables
    \item Raw JSON result files for independent verification
\end{itemize}

